\begin{frame}{실습: 가상 발표를 코드로 검증하기}
  8.2 실습(``99.8\%'' 검증)과 같은 구조로 가상 발표(``MLP로
  0.95'')를 코드로 검증 — 노트북이 \textbf{순서대로} 보여줌:
  \begin{enumerate}
    \item \textbf{세 모델 + 베이스라인}을 \textbf{같은 test
      분할}에서 정확도·F1·PR-AUC·AUC·log-loss로 비교 —
      ``0.95 vs 0.94''가 실제로 얼마나 큰 차이인지
    \item \textbf{MLP 학습 곡선} — 에폭 단위로 그려 ``에폭 50
      이 정당한가''를 val \textbf{포화 지점}으로 확인
    \item \textbf{표본 잡음 실험} — \textit{같은 모델·가중치}로
      test 100건을 \textbf{60번} 하위샘플링, 모델 간 차이가
      \textbf{표본 선택}에 의존하는지
    \item \textbf{교정 체크} — ``\textbf{0.90+ 확률}로 예측한
      건''이 실제로 얼마나 맞는지를 계수로 — 정확도 0.95가
      \textbf{신뢰할 수 있는 확률}인지 분리
  \end{enumerate}
  \vskip0.4em
  \begin{columns}[T]
    \begin{column}{0.4\textwidth}
      \centering
      \includegraphics[width=0.88\linewidth]{ch16_2_mlp_epoch_curve.png}
      \vskip0.3em
      {\footnotesize val 정확도 25~30 부근 포화, 그 뒤 train
      손실만 계속 하락}
    \end{column}
    \begin{column}{0.58\textwidth}
      {\footnotesize 노트북의 핵심 숫자 \textbf{세 가지}를 미리
      짚으면, ``이 숫자가 나오면 3점 질문을 쓸 수 있다''는
      감각이 생긴다. (다음 프레임)}
    \end{column}
  \end{columns}
\end{frame}
